\chapter{Ethical Considerations}
\label{app:ethics}

This appendix discusses ethical considerations for deploying ML in healthcare.

%==============================================================================
\section{Ethical Framework}
%==============================================================================

\begin{table}[h]
\centering
\caption{Ethical Considerations}
\begin{tabular}{ll}
\toprule
Issue & Approach \\
\midrule
Bias \& fairness & Stratified evaluation by age, sex; report per-node metrics. \\
Privacy & Federated learning; no raw data sharing; differential privacy option. \\
Transparency & SHAP explanations; case studies; interpretable baselines. \\
Accountability & Model as decision support; clinician retains final say. \\
Consent & MIMIC de-identified; institutional approval for research use. \\
\bottomrule
\end{tabular}
\label{tab:ethics_full}
\end{table}

%==============================================================================
\section{Policy Implications}
%==============================================================================

\textbf{GDPR-compliant deployment:} Federated learning aligns with data minimization and purpose limitation. Model updates (gradients) can be aggregated without sharing raw records. Local data remains at each site.

\textbf{Cost savings:} A 1\% accuracy gain (e.g., AUROC 0.76 vs.\ 0.75) could reduce unnecessary ICU days or inappropriate discharges. Economic analysis would require site-specific cost data.

%==============================================================================
\section{Patent-Style Claims}
%==============================================================================

For reference, a method claim could read:

\begin{quote}
\emph{``A method for few-shot token-based EHR prediction comprising: (a) extracting lab and vital measurements from a first 24-hour window of an ICU stay; (b) tokenizing said measurements into discrete symptom tokens with intensity weighting; (c) encoding said tokens via a transformer encoder; (d) computing class prototypes from support set embeddings; (e) predicting treatment feasibility for a query patient by distance to said prototypes.''}
\end{quote}
